\section{Normal Flow dan Positioning}

\textbf{Normal flow} adalah model tata letak default browser di mana elemen HTML disusun berdasarkan urutan kemunculannya di dokumen. Elemen block (seperti \code{div}, \code{p}, \code{h1}) ditumpuk secara vertikal---masing-masing memulai baris baru dan mengisi lebar penuh kontainer. Elemen inline (seperti \code{span}, \code{a}, \code{strong}) mengalir horizontal dalam baris teks. Memahami normal flow adalah prasyarat untuk memahami bagaimana positioning CSS memodifikasi perilaku ini \cite{mdnCss}.

Properti \code{position} mengontrol bagaimana elemen diposisikan relatif terhadap normal flow. CSS menyediakan lima nilai positioning: \code{static} (default), \code{relative}, \code{absolute}, \code{fixed}, dan \code{sticky}. Setiap nilai memiliki konteks referensi dan perilaku yang berbeda terhadap layout elemen lain. Properti pendamping \code{top}, \code{right}, \code{bottom}, dan \code{left} menentukan offset dari posisi referensi; properti ini hanya berlaku jika \code{position} bukan \code{static} \cite{webdevCss}.

\begin{table}[htbp]
\centering
\begin{tabular}{|P{2cm}|P{3cm}|P{3.5cm}|P{3.5cm}|}
\hline
\textbf{Nilai} & \textbf{Referensi} & \textbf{Efek pada Flow} & \textbf{Kasus Penggunaan} \\
\hline
\code{static} & Normal flow & Tetap di flow & Default; tidak perlu diset \\
\hline
\code{relative} & Posisi asal sendiri & Tetap di flow; ruang asli dipertahankan & Konteks untuk \code{absolute} child \\
\hline
\code{absolute} & Ancestor terdekat non-static & Keluar dari flow & Tooltip, dropdown, badge \\
\hline
\code{fixed} & Viewport & Keluar dari flow; tetap saat scroll & Header/navbar tetap \\
\hline
\code{sticky} & Viewport + scroll threshold & Di flow sampai threshold & Header sticky saat scroll \\
\hline
\end{tabular}
\caption{Lima nilai CSS position dan karakteristiknya}
\end{table}

Elemen \code{relative} tetap di posisi normal flow-nya, tetapi dapat digeser menggunakan \code{top}/\code{left}/\code{right}/\code{bottom} tanpa mempengaruhi elemen lain. Kegunaan utamanya adalah sebagai \textbf{containing block} untuk elemen \code{absolute} di dalamnya. Elemen \code{absolute} keluar sepenuhnya dari normal flow dan diposisikan relatif terhadap ancestor terdekat yang memiliki \code{position} non-static. Jika tidak ada ancestor seperti itu, elemen absolute merujuk ke viewport. Properti \code{z-index} mengontrol urutan tumpukan saat elemen overlap \cite{cssTricks}.

\begin{konsep}
\textbf{Stacking Context dan z-index.} Properti \code{z-index} hanya berlaku pada elemen yang memiliki \code{position} non-static (atau elemen flex/grid child). Nilai lebih tinggi ditempatkan di atas. Setiap elemen dengan \code{z-index} menciptakan \textit{stacking context} baru---elemen di dalam konteks tersebut tidak dapat ``menembus'' keluar. Ini menjelaskan mengapa kadang \code{z-index: 9999} pun tidak berhasil: elemen berada di stacking context yang lebih rendah dari elemen lain \cite{mdnCss}.
\end{konsep}

